
\subsubsection{6.1 Competing Explanations for the Observed Result}

Three explanations for Pythia's MMLU advantage are assessed.

\textbf{Architecture Difference (HIGH plausibility).} The most plausible explanation is the uncontrolled architecture difference. Pythia-6.9B uses GPT-NeoX architecture; OLMo-7B uses a LLaMA-style architecture with rotary positional embeddings and different attention configurations. Architecture-dependent differences in few-shot multiple-choice task performance have been observed across model families \citep{Touvron2023LLaMA}. An architecture advantage in MMLU would produce broadly distributed gains (consistent with Figure 3), would not affect HellaSwag (which is commonsense-based rather than knowledge-intensive), and would produce a stable gap across training checkpoints rather than a narrowing one. Without temporal trajectory analysis comparing both models at an earlier checkpoint (approximately 143 billion tokens), the stability versus narrowing of the gap cannot be determined, and causal attribution to corpus quality is not possible.

\textbf{Training Scale Insufficient (HIGH plausibility).} A second compelling explanation is that 300 billion training tokens is insufficient for Dolma's curation quality advantage to manifest. Groeneveld et al. [2024] report competitive MMLU performance for full-training OLMo-7B (approximately 2 trillion tokens), attributing part of this to Dolma's academic content inclusion. If the academic content advantage compounds with training scale --- as suggested by the data quality $\times$ token count interaction documented by Muennighoff et al. [2023] --- the Dolma advantage may not be detectable at 300 billion tokens.

\textbf{The Pile MMLU Contamination (LOW-MEDIUM plausibility).} The Pile may contain MMLU-adjacent content (e.g., academic PDFs from PubMed, ArXiv) that inflates Pythia's MMLU scores via benchmark contamination. However, the per-subject heatmap (Figure 3) shows Pythia's advantage is distributed across subjects less likely to appear in The Pile's specific domain sources (e.g., philosophy, sociology, moral scenarios), reducing the plausibility of contamination as the primary explanation.

\subsubsection{6.2 Structural Insight: HellaSwag Denominator Saturation}

The most methodologically significant finding is the HellaSwag convergence. Both models achieve exactly 0.4580 on HellaSwag at approximately 300 billion tokens, regardless of corpus quality. This convergence has a direct implication for metric design: the MMLU/HellaSwag ratio is not a valid discriminator of corpus curation quality at this training scale in a cross-architecture comparison, because the denominator provides no discriminative information.

A ratio metric $r = \text{numerator} / \text{denominator}$ is a reliable quality discriminator only when both components remain sensitive to the factor being tested. When the denominator saturates --- reaching a scale-dependent ceiling for the model family --- the ratio collapses into a scalar multiple of the numerator alone. For cross-architecture comparisons, numerator differences may reflect architecture effects rather than corpus quality. Researchers designing generalization balance metrics should verify denominator sensitivity at their target training scale before using ratio metrics as quality proxies.

\subsubsection{6.3 Limitations}

\textbf{L1: Architecture Confound Not Bounded.} The planned method for bounding the architecture confound --- temporal trajectory analysis comparing both models at approximately 143 billion and 300 billion tokens --- was not executed. The temporal trajectory would determine whether the ratio gap is widening (consistent with a quality-driven advantage that compounds with tokens) or stable (consistent with a fixed architecture-driven difference). Without this analysis, causal attribution of any observed performance difference to corpus quality is not possible. Future work should evaluate both models at an earlier checkpoint (Pythia step72000 $\approx$ 143 billion tokens; matched OLMo intermediate checkpoint) to distinguish these explanations.

\textbf{L2: Fast Evaluation (500-Sample Limit).} The use of \texttt{--limit 500} evaluates approximately 8--9 examples per MMLU subject, introducing high variance in per-subject accuracy estimates. Full evaluation across 14,042 MMLU questions would provide more reliable per-subject estimates and reduce bootstrap CI width. The large effect size (Cohen's d = $-2.732)$ and CI entirely below zero indicate that the directional refutation is unlikely to reverse with full evaluation, but the exact magnitude of $-0.0267$ is less reliable.

\textbf{L3: Single Training Scale.} Results are specific to approximately 300 billion training tokens. The corpus-quality hypothesis may hold at different scales; prior work suggests it does at approximately 2 trillion tokens \citep{Groeneveld2024OLMo}. The null result should not be generalized beyond its explicit scope.

\textbf{L4: Corpus Quality IV Not Directly Measured.} The corpus quality independent variable was operationalized via corpus identity (The Pile vs. Dolma) rather than direct measurement of quality proxy scores (n-gram repetition rate, Flesch-Kincaid grade level, fastText language identification confidence). Dolma's curation advantages are well-documented \citep{Soldaini2024Dolma}, but whether these advantages manifest in the specific quality dimensions expected to predict MMLU/HellaSwag balance has not been empirically verified on matched document samples.

\subsubsection{6.4 Implications for Metric Design}

Two design requirements for future generalization balance metrics follow from these results:

\begin{enumerate}
\item \textbf{Verify denominator sensitivity.} Before using a ratio metric as a quality proxy, confirm that the denominator task remains sensitive to the factor being tested at the target training scale. HellaSwag's convergence to identical values at approximately 300 billion tokens for both model families suggests it is an insensitive denominator for corpus quality comparisons at this scale.
\end{enumerate}

\begin{enumerate}
\item \textbf{Require architecture-matched designs for causal claims.} Single-checkpoint cross-architecture comparisons cannot isolate corpus quality effects. Future studies should either (a) use models with identical architecture trained on different corpora, or (b) use temporal trajectory analysis to bound the architecture confound within a cross-suite comparison.
\end{enumerate}

\subsubsection{6.5 Broader Context}

The primary positive contribution of this work is methodological: documenting where cross-architecture, single-checkpoint comparisons break down reduces the risk of misattributing performance differences to corpus quality when they may reflect architecture effects. A potential concern with publishing a null result on corpus curation quality is misinterpretation as evidence that curation does not matter in general. The results are explicitly scoped to approximately 300 billion training tokens, this specific architecture pair, and the MMLU/HellaSwag ratio metric. They do not contradict prior evidence of curation benefits at full training scale [Groeneveld et al., 2024; Penedo et al., 2023]; rather, they reveal the conditions under which those benefits are not detectable with current evaluation protocols.

